\section{Complete short-twig maps}\label{c:twig-appendix}

This appendix completes the local calculation used in
Lemma~\ref{c:twig}. Attach a path \(p-u-v\), where \(u\) has exactly
the two neighbours \(p,v\) and \(v\) is a leaf. Write \(de\) for the
directions \(p\to u\) and \(u\to v\). The condition \(e\ne-d\)
ensures that the three vertices are distinct and leaves exactly twelve
words. We stop upon returning to \(p\), so its other neighbours and
outgoing rule play no role in this table.

\begin{table}[htbp]
\centering\small
\setlength{\tabcolsep}{4pt}
\begin{tabular}{lllllll}
\toprule
\(de\) & Masks at \(u,v\) & \(u_0\to\) & \(u_1\to\)
& \(v_0\to\) & \(v_1\to\) & Return or internal cycle \\
\midrule
\(\N\N\)&\(\N\Sdir,\Sdir\)&\(p_1\)&\(p_1\)&\(u_1\)&\(u_1\)&All return as 1\\
\(\N\E\)&\(\E\Sdir,\W\)&\(p_1\)&\(v_1\)&\(u_1\)&\(u_0\)&All return as 1\\
\(\N\W\)&\(\Sdir\W,\E\)&\(v_0\)&\(p_1\)&\(u_1\)&\(u_1\)&All return as 1\\
\(\E\N\)&\(\N\W,\Sdir\)&\(v_1\)&\(p_0\)&\(u_1\)&\(u_1\)&All return as 0\\
\(\E\E\)&\(\E\W,\W\)&\(p_0\)&\(v_1\)&\(u_1\)&\(u_0\)&All return as 0\\
\(\E\Sdir\)&\(\Sdir\W,\N\)&\(p_0\)&\(v_1\)&\(u_0\)&\(u_1\)&\(u_0,v_0\) return; \((u_1,v_1)\)\\
\(\Sdir\E\)&\(\N\E,\W\)&\(p_0\)&\(v_1\)&\(u_1\)&\(u_0\)&All return as 0\\
\(\Sdir\Sdir\)&\(\N\Sdir,\N\)&\(v_1\)&\(v_1\)&\(u_0\)&\(u_1\)&All reach \((u_1,v_1)\)\\
\(\Sdir\W\)&\(\N\W,\E\)&\(p_1\)&\(v_0\)&\(u_1\)&\(u_1\)&\(u_0\) returns; \((u_1,v_0)\)\\
\(\W\N\)&\(\N\E,\Sdir\)&\(v_0\)&\(p_1\)&\(u_1\)&\(u_1\)&All return as 1\\
\(\W\Sdir\)&\(\E\Sdir,\N\)&\(v_1\)&\(p_1\)&\(u_0\)&\(u_1\)&All return as 1\\
\(\W\W\)&\(\E\W,\E\)&\(v_0\)&\(p_1\)&\(u_1\)&\(u_1\)&All return as 1\\
\bottomrule
\end{tabular}
\caption{Complete stopped maps for two-cell twigs. Every arrow is a
substitution in Table~\ref{tab:astar}. Parentheses in the final column
denote the unique internal tagged cycle in the specified row.}
\label{c:twig-table}
\end{table}

For the \(\N\E\) row, the longest return is
\(v_0\to u_1\to v_1\to u_0\to p_1\); all other starting tags give
suffixes of this route. The \(\E\E\) and \(\Sdir\E\) rows have the
same four-step form, returning to \(p_0\). In each of the remaining
nonabsorbing rows, the printed arrows give a return in at most four steps.
For \(\E\Sdir\), the tags \(u_1,v_1\) form the displayed cycle and
\(v_0\to u_0\to p_0\). For \(\Sdir\Sdir\),
\(v_0\to u_0\to v_1\) joins \((u_1,v_1)\). For \(\Sdir\W\),
\(v_1\to u_1\) joins \((u_1,v_0)\), while \(u_0\to p_1\).
These routes verify every transient and every internal recurrent component.

For a one-cell twig in direction \(d\), the leaf has mask \(\{-d\}\).
The endpoint rows therefore return state one for
\(d\in\{\N,\W\}\), preserve the arriving state for \(d=\Sdir\),
and complement it for \(d=\E\).

\subsection*{The permitted square ports}

Use the absolute square coordinates
\(a=(0,0),b=(1,0),c=(1,1),d=(0,1)\). In family \(\mathsf Q\)
there are at most two exterior cells and no second square. For a two-cell
twig at \(c\N\), a west continuation would meet the northwest corner,
so the only words are \(\N\N,\N\E\). The same grid-edge check gives
\(\E\E,\E\N\) at \(c\E\),
\(\N\N,\N\W\) at \(d\N\), and
\(\W\W,\W\N\) at \(d\W\).
All these rows return every tag. This excludes an attractor in a dormant
twig even for a trajectory that starts there.

The four active ports have the complete behaviour in
Table~\ref{c:active-ports}. The entry tag is obtained from the square
corner table. The one-cell column records the return of that entry tag;
an arbitrary tag on a one-cell twig also returns immediately, possibly in
the other state. The all-tag assertions in the two-cell column include
all four tags, not merely the tag entering from the core.

\begin{table}[htbp]
\centering\small
\begin{tabularx}{\textwidth}{llllX}
\toprule
Port & Entry tag & One cell & Two-cell word & Complete effect \\
\midrule
\(a\Sdir\)&\(u_1\)&Return \(a_1\)&\(\Sdir\Sdir\) or \(\Sdir\W\)&
The entering tag reaches the internal edge cycle. In the second word only
\(u_0\) returns to \(a_1\).\\
\addlinespace
\(a\W\)&\(u_0\)&Return \(a_1\)&\(\W\W\) or \(\W\Sdir\)&
Every tag returns to \(a_1\).\\
\addlinespace
\(b\E\)&\(u_1\)&Return \(b_0\)&\(\E\E\) or \(\E\Sdir\)&
The first word returns every tag to \(b_0\). The second captures the
entry tag; \(u_0,v_0\) return to \(b_0\).\\
\addlinespace
\(b\Sdir\)&\(u_0\)&Return \(b_0\)&\(\Sdir\E\) or \(\Sdir\Sdir\)&
The first word returns every tag to \(b_0\). The second captures every
tag on its outer state-one edge.\\
\bottomrule
\end{tabularx}
\caption{Active square-port returns and absorbers. The one-cell return
uses the listed entry tag. The two-cell options exclude continuations
that make a second square.}
\label{c:active-ports}
\end{table}

In the \(a\Sdir\) case with word \(\Sdir\W\), the returnable tag
does not become trapped within the stopped twig map: the square sends it
\(u_0\to a_1\to b_1\to a_0\to u_1\), after which it enters the
absorber. Likewise the returnable tags of the \(\E\Sdir\) twig at
\(b\E\) return to \(b_0\), travel through the square back to
\(b_1\), and are then sent to \(u_1\). These uses of the square core
are part of the proof of Lemma~\ref{c:square}, rather than extra assertions
about an isolated twig.
